\documentclass[11pt]{article}
\usepackage[a4paper,margin=25mm]{geometry}
\usepackage{amsmath,amssymb,amsthm,mathtools,booktabs,microtype}
\usepackage[T1]{fontenc}
\usepackage{lmodern}
\usepackage[colorlinks=true,linkcolor=blue,urlcolor=blue]{hyperref}
\newtheorem{theorem}{Theorem}[section]
\newtheorem{lemma}[theorem]{Lemma}
\newtheorem{corollary}[theorem]{Corollary}
\newtheorem{proposition}[theorem]{Proposition}
\newcommand{\F}{\mathbb F}
\newcommand{\SV}{\operatorname{SV}}
\newcommand{\rank}{\operatorname{rank}}
\newcommand{\ord}{\operatorname{ord}}
\newcommand{\GL}{\operatorname{GL}}
\newcommand{\Span}{\operatorname{span}}
\newcommand{\pdeg}{\operatorname{pdeg}}
\newcommand{\wt}{\operatorname{wt}}
\title{A Nine-Parameter Hitting Map for Affine Orbits\\
of Width-Three ROABPs in Every Characteristic}
\author{icekylinx \and GPT-6 Astra Pro}
\date{26 October 2026}
\begin{document}
\maketitle

\begin{abstract}
We give a self-contained argument for a polynomial-size black-box
hitting set for the invertible affine orbit of width-three read-once
oblivious algebraic branching programs, with unknown variable order,
over every field characteristic. The number of parameters is nine and
does not depend on the degree or the characteristic. The standard
five-pair Shpilka--Volkovich map also hits this class.

The proof uses a separating rational curve, a three-coordinate Hasse
jet lemma, and a rank-two Schur identity. After suitable constant
gauges, the whole first nonzero homogeneous scalar term has base-
characteristic digit weight at most two. Two fresh SV pairs detect
that term. A preliminary two-pair stage preserves transfer ranks and
boundary vectors; splitting rank-one transfers then reduces the
general word to the required mixed rank-two/invertible segments.
\end{abstract}

\noindent\textbf{Review status.}
This manuscript has received internal independent mathematical and
exact-source checks, in addition to the audits of its preceding proof
notes. It has not undergone external peer review, and no publication
priority claim is made.

\section{Model, map, and scope}

Let \(n,\delta\ge1\), let \(A_i(T)\in\F[T]^{3\times3}\) have entry
degrees at most \(\delta\), and let \(u,v\) be constant row and column
endpoints. Let \(L\in\GL_n(\F)\) have coordinate forms
\(\ell_1,\ldots,\ell_n\), and let \(c\in\F^n\). We consider
\begin{equation}\label{eq:model}
 f(X)=uA_1(\ell_1(X)+c_1)\cdots A_n(\ell_n(X)+c_n)v.
\end{equation}
The order, \(L\), \(c\), and the representation are unknown to the
black-box algorithm. Width at most three is included by zero padding;
unused coordinates and constant factors are allowed. The local generic
ranks may be one, two, or three. A zero local factor cannot occur in a
nonzero scalar word.

Fix distinct nodes \(a_1,\ldots,a_n\) in a working field, and let
\(\lambda_i(y)\) be their Lagrange interpolation polynomials. Write
\[
 \SV_{h,i}=\sum_{j=1}^h z_j\lambda_i(y_j).
\]
In characteristic zero put \(e_i=i\). In characteristic \(p>0\), let
\(e_1<\cdots<e_n\) be the first \(n\) positive integers not divisible
by \(p\). In both cases \(e_n\le2n-1\). Define
\begin{equation}\label{eq:map}
 \mathcal H_i(t,\boldsymbol y,\boldsymbol z)
   =t^{e_i}+\sum_{j=1}^4z_j\lambda_i(y_j).
\end{equation}

\begin{theorem}\label{thm:main}
For every field characteristic and every nonzero polynomial
\eqref{eq:model},
\[
 f(\mathcal H)\ne0.
\]
The standard map \(\SV_5\) also hits every such polynomial. A Cartesian
grid gives at most
\[
 \bigl((2n-1)n\delta+1\bigr)^9
\]
oracle evaluations for \eqref{eq:map}, or \((n^2\delta+1)^{10}\) for
\(\SV_5\). Extension-field evaluations are permitted when the base
field is too small. For \(n=0\) or \(\delta=0\), the model is constant
and one evaluation suffices.
\end{theorem}

The statement concerns an invertible affine orbit, not arbitrary singular
projections or an orbit closure. If a promise bounds only the individual
degrees of the latent scalar polynomial, each local matrix may be
truncated above that bound: its variable occurs in no other factor, so
the discarded terms contribute only scalar monomials exceeding the
promised individual degree.

For context, Bhargava--Ghosh~\cite{BG} give a quasipolynomial bound for
unknown-order ROABP affine orbits in characteristic zero or greater than
the latent individual degree. Their polynomial bound at constant width
for \emph{any-order} ROABPs uses the stronger promise that a width-bounded
representation exists in every variable order. The present statement
uses one unknown order and makes no restriction on the characteristic.
This comparison is not a certification that no later or equivalent
result exists.

\section{Three-coordinate rational Hasse jets}

For a rational vector \(R(t)\), let \(D^{(m)}R\) be the coefficient
of \(h^m\) in \(R(t+h)\). Call \(m\) a \emph{pivot order} if
\(D^{(m)}R\) is not in the \(\F(t)\)-span of its earlier Hasse rows.
In characteristic \(p\), write \(\wt_p(m)\) for the sum of the base-
\(p\) digits of \(m\).

\begin{lemma}\label{lem:jets}
Let \(0\ne R(t)\in\F(t)^q\), with \(q\le3\). In characteristic
zero every pivot order is at most \(q-1\). In characteristic \(p>0\)
every pivot order \(m\) satisfies \(\wt_p(m)\le q-1\).
More precisely, the first positive pivot, if present, is a \(p\)-power;
a third pivot, if present, is either twice that first power or another
\(p\)-power. In characteristic zero the only possible positive pivots
are one and two.
\end{lemma}

\begin{proof}
We may extend the field to an algebraic closure, since nonzero minors
remain nonzero. Divide the vector by a nonzero coordinate. The product
rule expresses each new Hasse row as a nonzero multiple of the old row
of that order plus earlier rows, so the pivot orders are unchanged.
It suffices to treat \((1,F_1(t),F_2(t))\), omitting missing coordinates.

The first nonzero Hasse variation of a nonconstant rational function in
characteristic \(p\) has order \(p^e\): extract its maximal Frobenius
power, after which its ordinary derivative is nonzero. In characteristic
zero that order is one. If both remaining functions are constant, there
is nothing more to prove.

In positive characteristic, take the largest common \(q_0=p^e\) for
which both functions belong to \(\F(t^{q_0})\), and put \(s=t^{q_0}\).
The pivot orders in \(t\) are \(q_0\) times those in \(s\). At least
one remaining function has nonzero derivative in \(s\); call it \(F\),
and call the other \(G\). In characteristic zero use \(s=t\).
The first positive pivot is now one. Put
\[
 h=G'/F',\qquad R_m=D^{(m)}G-hD^{(m)}F\quad(m\ge2).
\]
The third pivot is the first order with \(R_m\ne0\), if any.

If \(h'\ne0\), the characteristic is not two: in characteristic two
both second ordinary derivatives are zero, which forces \(h'=0\).
The product rule gives \(R_2=h'F'/2\ne0\). Thus the third pivot is two.

Suppose \(h'=0\). In characteristic zero, \(h\) is constant and
\(G-hF\) is constant, so there is no third pivot. In positive
characteristic both \(h\) and \(J=G-hF\) lie in \(\F(s^p)\).
If \(h\) is constant, the first nonzero Hasse variation of a nonconstant
\(J\) gives a \(p\)-power third pivot; if \(J\) is constant there is
none.

If \(h\) is nonconstant, let \(a=p^e\ge p\) be its first positive
Hasse order. Let \(b\) be the first positive Hasse order of \(J\), or
infinity if \(J\) is constant. For \(m<a\) we have \(R_m=D^{(m)}J\).
If \(b<a\), the third pivot is the \(p\)-power \(b\). If \(b\ge a\),
all these earlier defects vanish and
\[
 R_a=D^{(a)}J+F D^{(a)}h.
\]
Both Hasse derivatives on the right belong to \(\F(s^p)\), since
\(J,h\) do and \(p\mid a\). Thus
\(R_a'=F'D^{(a)}h\ne0\), proving that the third pivot is \(a\).
Multiplying the orders back by \(q_0\) proves the assertion. The
initial extension removes any need for perfectness of the original
coefficient field.
\end{proof}

This is a small-dimensional instance of the classical generic-order
criterion associated with Hasse derivatives; see
Stoehr--Voloch~\cite{SVorder}. The complete rational-function argument
above is the only version needed here. If a matrix \(Q(t)\) is applied
to \(R(t+h)\), it is always held fixed in \(h\). Multiplying every
Hasse row by that invertible matrix preserves pivot ranks. We never
substitute the derivatives of \(R(t)Q(t)\) for those of \(R(t+h)Q(t)\).

\section{Detecting a polynomial of digit degree two}

In positive characteristic define the \emph{digit degree} of a monomial
as the sum of the base-\(p\) digits of all its exponents. In
characteristic zero use its ordinary degree. The digit degree of a
polynomial is the largest such value among its nonzero monomials.

\begin{lemma}\label{lem:digit}
Over every coefficient-field extension, \(\SV_2\) hits every nonzero
homogeneous polynomial of digit degree at most two. An invertible
linear substitution preserves that assertion.
\end{lemma}

\begin{proof}
Linear substitution preserves digit degree because a \(p^e\)-th power
of a linear form is a linear combination of \(p^e\)-th powers. It
preserves nonzeroness when the linear map is invertible.

For each \(e\), the polynomials \(\lambda_i(y)^{p^e}\) are linearly
independent over every coefficient extension: evaluation at \(a_j\)
gives the identity matrix. Thus one pair detects a nonzero sum of
equal pure powers.

At degree \(p^e+p^f\) with \(e\ne f\), write the polynomial as
\(\sum_{i,j}c_{ij}X_i^{p^e}X_j^{p^f}\). The coefficient of
\(z_1^{p^e}z_2^{p^f}\) after two pairs is
\[
 \sum_{i,j}c_{ij}\lambda_i(y_1)^{p^e}\lambda_j(y_2)^{p^f},
\]
which is nonzero by independence in the two parameter variables.
At degree \(2p^e\) in odd characteristic, ordinary quadratic
polarization in \(X_i^{p^e}\) gives a nonzero bilinear coefficient.
In characteristic two at degree \(2p^e=p^{e+1}\), any cross term
with two distinct variables gives a nonzero symmetric bilinear
coefficient of \(z_1^{p^e}z_2^{p^e}\); pure powers do not contribute to
it. If there are no cross terms, the pure-power case applies.
Characteristic zero is the usual linear/quadratic argument. These
cases exhaust homogeneous digit degree at most two.
\end{proof}

\section{The lowest term of an identity/projector word}

Put \(P=\operatorname{diag}(1,1,0)\). Let \(B_i(Y_i)\) be rational
matrices over a field \(E\), regular at zero, on distinct variables,
with constants \(I_3\) or \(P\). A factor with constant \(P\) is
assumed to have generic rank two. Consider
\[
 p(Y)=e_1^{\mathsf T}B_1(Y_1)\cdots B_N(Y_N)e_2.
\]
Write
\[
 b_i=(B_i)_{13},\quad c_i=(B_i)_{32},\quad d_i=(B_i)_{12},
 \qquad
 \alpha_i=\ord b_i,\quad\beta_i=\ord c_i,\quad\delta_i=\ord d_i,
\]
where \(\ord0=\infty\). A pair \(i<j\) is \emph{allowed} when every
strict intermediate constant is \(I_3\). Its endpoints may have
constant \(P\).

\begin{lemma}\label{lem:lowest}
If \(p\ne0\), its first nonzero total degree is
\begin{equation}\label{eq:min}
 m=\min\left\{\min_i\delta_i,
         \min_{\text{allowed }i<j}(\alpha_i+\beta_j)\right\}.
\end{equation}
Its whole homogeneous term of degree \(m\) is
\begin{equation}\label{eq:leading}
 \sum_{\delta_i=m}d_{i,m}Y_i^m+
 \sum_{\substack{\text{allowed }i<j\\\alpha_i+\beta_j=m}}
 b_{i,\alpha_i}c_{j,\beta_j}Y_i^{\alpha_i}Y_j^{\beta_j},
\end{equation}
where the subscripts denote the corresponding nonzero coefficients.
\end{lemma}

\begin{proof}
At a \(P\)-layer let \(E_i\) be the active \(2\)-by-\(2\) block.
Its determinant is a unit. The rank-two Schur identity is
\begin{align}
 (B_i)_{33}&=b_iU_i+c_iV_i+d_iW_i,\label{eq:schur}\\
 U_i&=\frac{(B_i)_{31}(B_i)_{22}-c_i(B_i)_{21}}{\det E_i},\nonumber\\
 V_i&=\frac{(B_i)_{11}(B_i)_{23}}{\det E_i},\qquad
 W_i=-\frac{(B_i)_{31}(B_i)_{23}}{\det E_i}.\nonumber
\end{align}
Here \(U_i,V_i\) have positive order and \(W_i\) has order at least
two. Thus a finite order of \((B_i)_{33}\) is strictly larger than
at least one of \(\alpha_i,\beta_i,\delta_i\).

Expand a coefficient path and examine its first visit to state two.
A direct transition from one to two costs at least \(\min\delta_i\).
Otherwise choose its last transition from one to three before that
visit; until the transition from three to two, the path stays in
state three. Without an intermediate \(P\)-layer these endpoints
give an allowed pair of no greater degree. At an intermediate
\(P\)-layer, \eqref{eq:schur} shortens the path: a \(b_i\) term permits
discarding the earlier inactive prefix, a \(c_i\) term permits
discarding the later inactive suffix, and a \(d_i\) term gives a direct
transition. The discarded multiplier has positive degree. Equivalently,
use the strict order inequality and recurse on a smaller interval.
Continue until a direct transition or an allowed pair is reached.
Every path crossing an intermediate \(P\)-layer therefore has degree
strictly larger than such a witness.

No coefficient can occur below \eqref{eq:min}. At that degree only
the direct steps or the allowed two-step paths remain. Their
coefficients do not cancel: a direct monomial has support one, and a
two-variable monomial has the unique intermediate state three. A path
using active state one or two instead would require a direct \(d\)
and an additional positive degree. Different pairs have different
supports. This proves \eqref{eq:leading}. If the minimum were infinite,
the shortening argument would make every scalar coefficient zero.
\end{proof}

\begin{lemma}[Low-degree reset]\label{lem:reset}
Let \(m\) be as above. At a \(P\)-layer the inactive coordinate of a
prefix row is, modulo terms of degree at least \(m\), its local
\(b_i\) multiplied by a unit with constant one. The transpose/reverse
statement holds for a suffix column and \(c_i\).
\end{lemma}

\begin{proof}
Write the row just before the layer as \((r_1,r_2,r_3)\), with constant
\(e_1^{\mathsf T}\). Every prefix's active second coordinate has order
at least \(m\), since later constant factors preserve it. The update
of that coordinate, together with \(\ord d_i\ge m\), gives
\(\ord(r_3c_i)\ge m\). Using \eqref{eq:schur}, the inactive update is
\[
 r_3'=(r_1+r_3U_i)b_i+r_2(B_i)_{23}
                 +(r_3c_i)V_i+(r_3d_i)W_i.
\]
The last three terms have order at least \(m\). The first multiplier
is a unit with constant one. Transpose and reverse for the dual result.
\end{proof}

\section{Constant gauges and actual rational curves}

A \emph{separating curve} here means a polynomial curve
\(\gamma(t)\in\F[t]^n\) whose coordinate derivatives are
\(\F\)-linearly independent. The curves in \eqref{eq:map} have this
property: their nonzero derivatives are distinct monomials.
For independent latent forms, every
\[
 C_i(t)=c_i+\ell_i(\gamma(t))
\]
has \(C_i'\ne0\), and hence is transcendental over the coefficient
field. The same remains true after adjoining independent old seeds
and adding their values to the constants \(c_i\).

\begin{lemma}[Pure active endpoint gauges]\label{lem:gauges}
Let \(F_i(T_i)\) be rational univariate matrices over a field \(\F_0\),
and evaluate them at \(C_i(t)\), with \(C_i'\ne0\). Assume local ranks
are two or three and are preserved by this evaluation. Assume every
constant interval containing a rank-two layer has rank two. Put
\(A_i=F_i(C_i)\), \(M=A_1\cdots A_N\), and take fixed endpoints \(u,v\)
such that
\[
 uM\ne0,\qquad Mv\ne0,\qquad uMv=0.
\]
There are rational cut gauges \(Q_i(t)\) for which
\[
 Q_{i-1}A_iQ_i^{-1}\in\{I_3,P\},\qquad
 u=e_1^{\mathsf T}Q_0,\qquad Q_Nv=e_2.
\]
The conclusion also holds when all factors are invertible.
\end{lemma}

\begin{proof}
Choose a row \(b\) with \(bMv=1\). The rows \(uM,bM\) are independent,
so \(U_0=\Span(u,b)\) maps injectively under \(M\). Propagate
\(U_i=U_0A_1\cdots A_i\) with these propagated basis rows.
They remain two-dimensional. Choose a one-dimensional final complement
\(W_N\) to \(U_N\) contained in \(\ker(v)\); it exists because
\(U_Nv\ne0\). Put
\[
 W_i=\{x:xA_{i+1}\cdots A_N\in W_N\}.
\]
When the tail contains a rank-two layer, it has rank two and image
\(U_N\), so \(W_i\) is its one-dimensional kernel. Otherwise the tail
is invertible and \(W_i\) is the one-dimensional inverse image of
\(W_N\). Thus \(U_i\oplus W_i\) is the full row space and
\(W_{i-1}A_i\subseteq W_i\). A rank-two layer kills the passive line;
an invertible layer maps it isomorphically. Normalize its basis along
each invertible run to make every nonzero passive map equal to one.
Let \(Q_i\) have the resulting three basis rows. The active maps are
identities, giving the asserted constants. The endpoint identities
follow from the first basis row \(u\), from \(uMv=0\), and from
\(bMv=1\). All choices use linear algebra over \(\F_0(t)\).
\end{proof}

For residual variables set
\[
 B_i(Y_i)=Q_{i-1}(t)F_i(C_i(t)+Y_i)Q_i(t)^{-1}.
\]
Let \(\Delta_i=C_i(t+h)-C_i(t)=C_i'(t)h+O(h^2)\).
The normalized prefix row after this substitution is
\begin{equation}\label{eq:prefix}
 \bigl[uF_1(C_1(t+h))\cdots F_j(C_j(t+h))\bigr]Q_j(t)^{-1}.
\end{equation}
The suffix column is \(Q_{j-1}(t)\) times its actual suffix curve at
\(t+h\). These are rational curves over \(\F_0\), followed or preceded
by invertible gauges held fixed in \(h\). No singular prefix matrix is
inverted.

\begin{lemma}[Mixed segment detection]\label{lem:mixed}
Under the rank and boundary hypotheses of Lemma~\ref{lem:gauges}, a
nonzero shifted scalar is detected by two fresh original-coordinate
SV pairs, when its residual variables are independent linear forms
\(Y_i=\ell_i(X)\). The same conclusion holds for a fully invertible
rational word with nonzero endpoint vectors.
\end{lemma}

\begin{proof}
If its value at \(C\) is nonzero, set the fresh amplitudes to zero.
Otherwise use the gauges and Lemma~\ref{lem:lowest}. Consider a minimal
allowed pair \((i,j)\), with \(\alpha_i+\beta_j=m\). Every \(b_h\)
in the identity run before \(j\), including its preceding \(P\)-layer
if present, has order at least \(\alpha_i\); otherwise \((h,j)\)
would be cheaper. Choose the first one attaining that order.
Lemma~\ref{lem:reset} discards older inactive terms below \(m\), and
the intervening identity constants show that this prefix's first
projective variation has exactly that order in its inactive coordinate.
Its coefficient after \(Y=\Delta\) is
\(b_{h,\alpha_i}(C_h')^{\alpha_i}\ne0\). Its active second coordinate
starts only at degree \(m>\alpha_i\). By \eqref{eq:prefix} and
Lemma~\ref{lem:jets}, \(\alpha_i\) is a \(p\)-power, or one in
characteristic zero. The last minimum \(c_h\) in the dual suffix run
similarly proves that \(\beta_j\) is a \(p\)-power. This applies to
every minimal pair.

If any such pair exists, \(m\) is the sum of two \(p\)-powers. Thus
all mixed terms in \eqref{eq:leading}, and all direct terms \(Y_i^m\),
have digit degree at most two. If there is no minimal pair, choose the
first direct factor with \(\delta_j=m\). Its prefix's active second
coordinate has the unique degree-\(m\) term \(d_{j,m}Y_j^m\).
After \(Y=\Delta\), its coefficient at \(h^m\) is nonzero, while all
earlier Hasse rows have zero active second coordinate. Therefore \(m\)
is a pivot of the actual three-coordinate curve in \eqref{eq:prefix}.
Lemma~\ref{lem:jets} again bounds its digit weight by two.

Consequently the \emph{whole} first nonzero homogeneous scalar term
has digit degree at most two. The independent residual forms give an
injective substitution \(Y_i=\ell_i(X)\), which preserves digit degree.
If the segment uses only a subset of the latent coordinates, extend
its independent forms to a full basis for this argument.
Lemma~\ref{lem:digit} detects the resulting nonzero term. Higher terms
have larger total fresh-amplitude degree and cannot cancel it. All
rational denominators are units at fresh amplitude zero.
\end{proof}

\section{The invertible precursor and endpoint completion}

A nonzero rational univariate row can be completed to a generically
invertible \(3\)-by-\(3\) matrix by using two \emph{constant} standard
rows over the original coefficient field. Choose a nonzero coordinate
of the row and use the other two standard rows. Do the column version
for a column endpoint. These choices are made before evaluating the
curve. Their determinants are nonzero univariate rational functions,
so remain nonzero at the transcendental \(C_i\). Constant endpoints
admit constant completions.

\begin{proposition}[Restricted precursor]\label{prop:precursor}
The map \(\gamma+\SV_2\) hits every nonzero scalar consisting of a
generically invertible width-three rational run and rational univariate
endpoint vectors on distinct latent variables. It also hits words whose
local generic ranks are one or three.
\end{proposition}

\begin{proof}
Complete the endpoints as just described. Every constant matrix after
the separating curve is invertible. Nonzero source and sink endpoints
therefore have nonzero active projections, and Lemma~\ref{lem:mixed}
applies in its all-invertible case.

For a rank-one factor \(R(T)\), choose a nonzero entry \(h=R_{ab}\)
and use the exact factorization
\[
 R(T)=(R(T)e_b)\bigl(e_a^{\mathsf T}R(T)/h(T)\bigr).
\]
Every pivot is nonzero at \(C_i\). Splitting at all such factors gives
a product of nonzero rational scalar segments with invertible middle
runs. The same fresh pairs hit every segment, so their product remains
nonzero. Neighboring scalar factors may share a boundary variable;
within each segment the variables remain distinct.
\end{proof}

\section{Rank preparation for a general width-three word}

Begin with the symbolic shift
\[
 G=\gamma(t)+\SV_2(\text{old seeds}),\qquad
 C_i=c_i+\ell_i(G).
\]
Split original rank-one factors by their univariate pivot
factorizations. Each resulting scalar chain has local ranks two or
three and univariate endpoint vectors. It is nonzero because their
product is nonzero. Chains with no rank-two factor are covered by
Proposition~\ref{prop:precursor}.

In a remaining chain, write its rank-two factors, in their order, as
\[
 R_i(T)=U_i(T)V_i(T),\qquad 1\le i\le k,
\]
where \(U_i\) is \(3\)-by-\(2\) of full column rank and \(V_i\) is
\(2\)-by-\(3\) of full row rank. Choose two independent columns and
invert a nonzero \(2\)-by-\(2\) pivot minor to obtain this rational
univariate factorization. Every denominator remains nonzero at \(C_i\).
Let \(M_i\) be the invertible run between \(R_i\) and \(R_{i+1}\),
with \(M_0,M_k\) the outer runs, and let \(l,r\) be the endpoints.
Define
\begin{equation}\label{eq:transfers}
 \alpha=lM_0U_1,\qquad \beta=V_kM_kr,\qquad
 T_i=V_iM_iU_{i+1}\quad(1\le i<k).
\end{equation}
The scalar is \(\alpha T_1\cdots T_{k-1}\beta\).

Every nonzero entry of \(\alpha,\beta,T_i\) is an invertible-run scalar
with univariate endpoint vectors. Proposition~\ref{prop:precursor}
therefore preserves it at \(G\), in particular
\(\alpha(C),\beta(C)\ne0\). A nonzero determinant of \(T_i\) is the
scalar
\[
 (\textstyle\bigwedge^2 V_i)
 (\textstyle\bigwedge^2 M_i)
 (\textstyle\bigwedge^2 U_{i+1}),
\]
whose middle word again consists of invertible \(3\)-by-\(3\) compound
matrices and whose endpoints are rational univariate vectors. It too
is preserved. Thus every transfer retains its generic rank, one or
two. Rank zero would make the original scalar chain zero.

All assertions hold simultaneously over the old-seed function field.
In particular, boundary nonzeroness has been proved separately; it is
not inferred from preservation of square interval ranks.

\section{Rank-one transfer splitting and original-variable recovery}

For each rank-one \(T_i\), choose a pivot
\(h_i=(T_i)_{a_i b_i}\ne0\), whose value at \(C\) is nonzero, and
factor
\[
 T_i=(T_ie_{b_i})\bigl(e_{a_i}^{\mathsf T}T_i/h_i\bigr).
\]
The pivot may depend on several variables, but is a scalar unit at
the shifted origin. Splitting at these transfers expresses the scalar
chain as a product of nonzero rational scalars.

For consecutive marked transfers \(i<j\), multiply the intervening
scalar by \(h_i\). Its numerator is exactly
\begin{equation}\label{eq:recombine}
 e_{a_i}^{\mathsf T}V_iM_i
 R_{i+1}M_{i+1}\cdots M_{j-1}R_j
 M_jU_{j+1}e_{b_j}.
\end{equation}
Every \(U_hV_h\) has recombined to the original univariate matrix
\(R_h\). Therefore this is a word on distinct original latent variables,
with univariate endpoints \(e_{a_i}^{\mathsf T}V_i\) and
\(U_{j+1}e_{b_j}\). The scalar denominator is kept outside this word.
Different scalar factors may share variables; each numerator remains
read once internally.

All internal transfers \(T_{i+1},\ldots,T_{j-1}\) are invertible at
\(C\). Every middle interval containing a rank-two original factor
therefore has constant rank two. The left boundary projection through
the whole constant middle is nonzero: it is the nonzero row
\(e_{a_i}^{\mathsf T}T_i\), followed by invertible two-state transfers,
a full-row-rank \(V_j\), and an invertible outer run. The right
projection is nonzero by the dual argument using \(T_je_{b_j}\).

For the first scalar factor, use the already nonzero row \(\alpha(C)\)
instead of a marked left transfer. For the last use the already nonzero
column \(\beta(C)\). If no transfer is marked, use both directly.
Adjacent marked transfers cause no exception: \eqref{eq:recombine}
then contains the single rank-two original factor \(R_{i+1}\).
These statements verify every boundary and interval hypothesis of
Lemma~\ref{lem:gauges} for each numerator after fixed univariate endpoint
completion.

\section{Completion of the proof and quantitative bounds}

\begin{proof}[Proof of Theorem~\ref{thm:main}]
Treat the old two SV pairs as constants and put
\(\F_0=\F(\text{old seeds})\). The parameter \(t\) is still
transcendental. Each \(C_i(t)\) still has nonzero derivative, and all
actual local factors and fixed endpoint completions are rational
univariate functions over \(\F_0\). The possibly imperfect nature of
\(\F_0\) is allowed by Lemma~\ref{lem:jets}.

Apply Lemma~\ref{lem:mixed} to every numerator from the preceding
section, using two new SV pairs. Every scalar pivot denominator stays
nonzero because its value at new amplitude zero was already nonzero.
A product of nonzero rational functions is nonzero. The exact
factorizations thus prove
\[
 f\bigl(\gamma+\SV_2(\text{old})+\SV_2(\text{new})\bigr)\ne0,
\]
which is \eqref{eq:map}.

For pure \(\SV_5\), specialize its first pair in characteristic \(p\)
to \(z_0=t,y_0=t^p\). The curve
\(\gamma_i=t\lambda_i(t^p)\) has independent derivatives
\(\lambda_i(t^p)\). In characteristic zero use \(z_0=t,y_0=t\);
the map \(q(t)\mapsto q(t)+tq'(t)\) is invertible on polynomials of
degree less than \(n\), so these derivatives are also independent.
The other four pairs remain fresh. Nonvanishing at this proof-only
specialization implies nonvanishing for unrestricted \(\SV_5\).

The original scalar has total degree at most \(n\delta\), unchanged
by an affine substitution. For the explicit curve in \eqref{eq:map},
the separate seed-degree bounds are
\[
 \deg_t\le e_n n\delta,\qquad
 \deg_{z_j}\le n\delta,\qquad
 \deg_{y_j}\le n(n-1)\delta.
\]
Consequently a more precise rectangular grid has size
\[
 (e_n n\delta+1)(n\delta+1)^4
                  \bigl(n(n-1)\delta+1\bigr)^4.
\]
The uniform bound in the theorem follows from \(e_n\le2n-1\).
For pure \(\SV_5\), every individual seed degree is at most
\(n^2\delta\), giving the stated ten-seed grid. The nonvanishing
statement concerns the original polynomial; numerical grid points
need not separately avoid every auxiliary rational pivot used in
the proof.
\end{proof}

\section{What the proof does and does not use}

The proof has no dependency on a general-width cone-preservation theorem,
on a coefficient-transport claim based only on digit weight, or on an
assumption that a generic rank-one transfer has univariate factors.
The low digit-degree assertion applies to the entire first homogeneous
scalar component after the actual separating-curve shift. Rank-one
transfer splitting is followed by exact recombination of the original
univariate matrices.

The decisive width-three feature is one inactive state beside a
two-dimensional active core. Its rank-two Schur entry lies in the ideal
of the two inactive boundary channels and the single active scalar
channel, with positive-order multipliers. No analogous general-width
statement is inferred here. Nor is the theorem a hitting result for
arbitrary singular affine projections, which can have very different
expressive power.

\begin{thebibliography}{9}
\bibitem{BG}
V. Bhargava and S. Ghosh.
\newblock Improved Hitting Set for Orbit of ROABPs.
\newblock \emph{Computational Complexity}, 2022.
\newblock \href{https://doi.org/10.1007/s00037-022-00230-9}{doi:10.1007/s00037-022-00230-9}.
\newblock Primary preprint:
\url{https://eccc.weizmann.ac.il/report/2021/062/revision/1/download/}.

\bibitem{SVorder}
K.-O. Stoehr and J. F. Voloch.
\newblock Weierstrass Points and Curves Over Finite Fields.
\newblock \emph{Proceedings of the London Mathematical Society},
52(1):1--19, 1986.
\newblock \href{https://doi.org/10.1112/plms/s3-52.1.1}{doi:10.1112/plms/s3-52.1.1}.
\end{thebibliography}
\end{document}
